\documentclass[12pt]{article}
\usepackage[T1]{fontenc}
\usepackage[utf8]{inputenc}
\usepackage{lmodern}
\usepackage{textcomp}
\usepackage{enumitem}
\usepackage{geometry}
\geometry{margin=1in}
\begin{document}
\begin{center}
Answer Key - Sociology - Undergraduate Level Assessment 21 \
\small (Answers are ordered exactly as in the question paper.)
\end{center}

\section*{Multiple Choice}
\begin{enumerate}[label=\arabic*.]
\item \textbf{Answer:} B
\\ \textit{Options:} A) Westernization theory; B) Modernization theory; C) Industrialization theory; D) Dependency theory
\\ \textit{Note:} Modernization theory posits that societies evolve through a linear process from traditional, agrarian structures to modern, industrialized ones, making it the appropriate term for this concept.
\item \textbf{Answer:} A
\\ \textit{Options:} A) those in non-manual occupations, exercising authority on behalf of the state; B) people working in consultancy firms who were recruited by big businesses; C) the young men and women employed in domestic service in the nineteenth century; D) those who had worked in the armed services
\\ \textit{Note:} Goldthorpe's concept of the 'service class' refers to individuals engaged in non-manual occupations who are responsible for administering and managing services, often within the framework of state institutions, highlighting their role in exercising authority and influence in society.
\item \textbf{Answer:} A
\\ \textit{Options:} A) we internalise and take on social roles from a pre-existing framework; B) we create and negotiate our roles through interaction with others; C) social roles are not fixed or stable but fluid and pluralistic; D) roles have to be learned to suppress unconscious motivations
\\ \textit{Note:} Role-learning theory posits that individuals acquire and embody social roles through a structured process of socialization within established frameworks of norms and expectations, thereby integrating these roles into their identity.
\item \textbf{Answer:} C
\\ \textit{Options:} A) the findings are amenable to statistical analysis; B) it is conducted over a period of several years; C) it uncovers rich, detailed accounts from an insider's perspective; D) it compares findings from a number of different cases
\\ \textit{Note:} Ethnographic research is designed to immerse the researcher in the community being studied, allowing for the collection of in-depth, nuanced descriptions and interpretations of social phenomena from the perspective of those within that culture, which ultimately results in qualitative data.
\item \textbf{Answer:} C
\\ \textit{Options:} A) genetics; B) evolution; C) height; D) brain size
\\ \textit{Note:} Height was not considered in the naturalistic explanations of race in the nineteenth century, as the focus was primarily on more prominent physical traits such as skull shape, skin color, and other visible characteristics that were mistakenly linked to perceived racial hierarchies.
\end{enumerate}

\section*{True / False}
\begin{enumerate}[label=\arabic*.]
\setcounter{enumi}{5}
\item \textbf{Answer:} False
\\ \textit{Options:} A) True; B) False
\\ \textit{Note:} Cultural restructuring often encompasses a broader approach that includes not only the preservation of historical sites but also the integration of economic revitalization efforts to enhance community development and sustainability.
\item \textbf{Answer:} True
\\ \textit{Options:} A) True; B) False
\\ \textit{Note:} Sociology is considered a social science because it employs systematic methodologies to formulate theories that are based on empirical research and can be validated or refuted through peer review, thereby ensuring the rigor and reliability of its findings.
\item \textbf{Answer:} False
\\ \textit{Options:} A) True; B) False
\\ \textit{Note:} The 'absolute' poverty line represents a minimum threshold of income necessary to meet basic needs for survival, but it does not necessarily capture all extreme levels of poverty that may exist within a society, particularly those that are relative or situational.
\item \textbf{Answer:} True
\\ \textit{Options:} A) True; B) False
\\ \textit{Note:} The post-war welfare state of 1948 primarily focused on social security, healthcare, and housing rather than establishing a minimum wage, which was not a central component of its objectives.
\item \textbf{Answer:} True
\\ \textit{Options:} A) True; B) False
\\ \textit{Note:} Incorporation in the context of the nineteenth-century labour movement refers to the process by which working-class organizations, such as trade unions, were recognized and included in political negotiations, thereby enabling them to advocate for workers' rights and interests within the political system.
\end{enumerate}

\section*{Long Form Response}
\begin{enumerate}[label=\arabic*.]
\setcounter{enumi}{10}
\item \textbf{Answer:} A population pyramid serves as a visual representation of the age and sex distribution within a population and can reveal significant demographic trends. In a population with a sustained higher birth rate than death rate, the pyramid would typically be wider at the bottom than at the top. This shape indicates a larger number of children and younger individuals compared to the older age groups, reflecting a robust birth rate that results in a youthful population. Over generations, this trend leads to a bulging base as more individuals enter the reproductive age, perpetuating the cycle of higher birth rates. Consequently, the population structure not only highlights immediate demographic trends but also hints at future growth potential and societal challenges related to age distribution.
\\ \textit{Note:} correct because a population pyramid with a wider base indicates a higher number of births relative to deaths, demonstrating a youthful demographic structure consistent with sustained higher birth rates over generations.
\item \textbf{Answer:} Official statistics on strike action in sociology are limited by the fact that many strikes go unnoticed by employers and the mass media, leading to significant underreporting. This lack of visibility can distort the understanding of workers' grievances and the overall landscape of labor struggles. When strikes are underreported, it diminishes their potential impact on social movements, as these actions often serve as catalysts for broader collective mobilization. Consequently, the failure to capture the full extent of strike activity can obscure the realities of workers' conditions and undermine the effectiveness of advocacy efforts aimed at labor reform. This highlights the need for more comprehensive data collection methods to accurately reflect the dynamics of strike action and its implications on societal change.
\\ \textit{Note:} correct because it highlights how underreported strikes obscure workers' grievances and weaken their potential to drive significant social movements, thereby distorting the overall understanding of labor struggles in sociology.
\end{enumerate}

\vfill
\noindent\textit{End of Answer Key}
\end{document}